\section{Radiation}\label{sec:radiation}

{\bf \Large
\begin{tabular}{ccc}
\hline
  Corresponding author & : & Yuta Kawai\\
\hline
\end{tabular}
}
\\

Radiative transfer describes the propagation of electromagnetic radiation through the atmosphere and determines the radiative heating and cooling resulting from the absorption, emission, and, in general, scattering of radiation.

\subsection{Basic formulation of radiative transfer}
%%
The following formulation follows the standard theory of radiative transfer~\citep[e.g.,][]{Liou2002}.
In general, the specific intensity depends on the position, propagation direction, wavelength, and time.
For a given wavelength and propagation direction, the radiative transfer equation is written as
%%
\begin{equation}
  \frac{dI_\lambda(\bm{\sigma})}{ds} =
  -\chi_\lambda I_\lambda(\bm{\sigma})
  +\kappa_\lambda B_\lambda
  +\sigma_{s,\lambda}\mathcal{S}_\lambda(\bm{\sigma}),
\label{eq:radiation_transfer_equation_general_form}
\end{equation}
%%
where $I_\lambda(\bm{\sigma})$ is the specific intensity at wavelength $\lambda$ in the propagation direction $\bm{\sigma}$,
$s$ is the distance along the radiation path,
$\kappa_\lambda$ and $\sigma_{s,\lambda}$ are the absorption and scattering coefficients, respectively,
$\chi_\lambda=\kappa_\lambda+\sigma_{s,\lambda}$ is the extinction coefficient,
and $B_\lambda$ is the Planck function. $\mathcal{S}_\lambda$ is the scattering source function defined as
%%
\begin{equation}
  \mathcal{S}_\lambda(\bm{\Omega}) = \frac{1}{4\pi}\int_{4\pi} I_\lambda(\bm{\Omega}') P_\lambda(\bm{\Omega},\bm{\Omega}') d\bm{\Omega}',
\end{equation}
%%
where $P_\lambda(\bm{\Omega},\bm{\Omega}')$ is the scattering phase function that describes the angular distribution of scattered radiation.
Introducing the optical thickness $\tau_\lambda$ defined as
%%
\begin{equation}
  d\tau_\lambda = \chi_\lambda ds,
\end{equation}
%%
Eq.\,\eqref{eq:radiation_transfer_equation_general_form} becomes
%%
\begin{equation}
  \frac{dI_\lambda}{d\tau_\lambda} = - I_\lambda + S_\lambda.
\label{eq:radiation_transfer_equation_general_form_tau}
\end{equation}
Here, $S_\lambda$ is the source function defined as
%%
\begin{align}
  S_\lambda &= \frac{\kappa_\lambda B_\lambda + \sigma_{s,\lambda} J_\lambda}{\chi_\lambda} \\
            &= (1-\omega_\lambda) B_\lambda + \omega_\lambda J_\lambda,
\end{align}
%%
where $\omega_\lambda = \sigma_{s,\lambda}/\chi_\lambda$ is the single scattering albedo.


In atmospheric simulations, the radiative transfer is typically solved independently in each atmospheric column under the plane-parallel approximation, 
whereby horizontal radiative transport is neglected.
We introduce the optical thickness along the vertical direction $z$ as
%%
\begin{equation}
  d\tau_{\lambda,z}= -\chi_\lambda dz= - \chi_\lambda \mu\, ds = -\mu\, d\tau_\lambda,
\end{equation}
%%
where $\mu$ is the cosine of the zenith angle $\theta$. 
Thus, we can rewrite Eq.\,\eqref{eq:radiation_transfer_equation_general_form_tau} as
%%
\begin{equation}
  \mu \frac{dI_\lambda}{d\tau_{\lambda,z}} = I_\lambda - S_\lambda. 
\label{eq:radiation_transfer_equation_general_form_plane_parallel}
\end{equation}
The equation for the radiative fluxes is obtained by integrating Eq.\,\eqref{eq:radiation_transfer_equation_general_form_plane_parallel} over the semi-sphere as
%%
\begin{align}
 \dfrac{dF_\lambda^\uparrow}{d\tau_{\lambda,z}} =&  
   4\pi J_\lambda^\uparrow 
 - 2\pi (1-\omega_\lambda) B_\lambda
 - \dfrac{\omega_\lambda}{2}\int_0^1 d\mu \int_{4\pi} P_\lambda(\bm{\Omega},\bm{\Omega}') I_\lambda(\bm{\Omega}') d\Omega', \\
  \dfrac{dF_\lambda^\downarrow}{d\tau_{\lambda,z}} =& 
 - 4\pi J_\lambda^\downarrow 
 + 2\pi (1-\omega_\lambda) B_\lambda
 + \dfrac{\omega_\lambda}{2}\int_{-1}^0 d\mu \int_{4\pi} P_\lambda(\bm{\Omega},\bm{\Omega}') I_\lambda(\bm{\Omega}') d\Omega'. \label{eq:rte_general_form_plane_parallel_flux}
\end{align}
%%
Here, superscripts $\uparrow$ and $\downarrow$ indicate the quantities integrated over the upward and downward hemispheres, respectively, 
$J_\lambda$ is the mean intensity defined as 
\begin{align}
  J_\lambda = \frac{1}{4\pi}\int_{4\pi} I_\lambda(\bm{\Omega}) d\Omega,
\end{align}
%%
and the upward and downward fluxes are defined as
%%
\begin{align}
F_\lambda^\uparrow =2\pi \int_0^1 I_\lambda(\mu)\mu\,d\mu,\;\;
F_\lambda^\downarrow =2\pi \int_{-1}^0 I_\lambda(\mu)|\mu|\,d\mu,
\end{align}
%%
where both fluxes are defined as positive quantities.



The hemispherically integrated equations are not closed because they still contain the angular moments of the radiation field and the scattering integral. 
In the two-stream approximation~\citep{Toon1989}, the angular dependence within each hemisphere is represented by a single effective propagation direction. 
The resulting angular integration is approximated using a diffusivity factor $D_\lambda$,
which accounts for the increased optical path length of obliquely propagating radiation. 
Then, the two-stream equations can therefore be written schematically as
\begin{align}
\frac{dF_\lambda^\uparrow}{d\tau_{\lambda,z}}
&= D_\lambda^\uparrow \left( F_\lambda^\uparrow -\pi S_\lambda^\uparrow \right), \\
\frac{dF_\lambda^\downarrow}{d\tau_{\lambda,z}}
&=
-D_\lambda^\downarrow \left( F_\lambda^\downarrow -\pi S_\lambda^\downarrow \right),
\label{eq:rte_twostream_approx}
\end{align}
%%
where $D_\lambda^\uparrow$ and $D_\lambda^\downarrow$ are the diffusivity factors determined by the adopted angular quadrature,
$S_\lambda^\uparrow$ and $S_\lambda^\downarrow$ denote the corresponding hemispherically averaged source functions.
Note that the source term $\pi S_\lambda$ is not obtained by directly evaluating $2\pi\int_0^1 S_\lambda\,d\mu$. 
Rather, it appears as part of the angular closure introduced in the two-stream approximation.


In the current version of SCALE-DG, 
the spectral dependence is treated separately in each radiation scheme by dividing the spectrum into one or more bands.
Furthermore, the scattering is neglected in all radiation schemes.
Consequently, 
$\sigma_{s,\lambda} = 0$ and $\chi_\lambda = \kappa_\lambda$, and the source function is simplified as $S_\lambda = B_\lambda$.
The radiation transfer equation in Eq.\,\eqref{eq:rte_twostream_approx} using the two-stream approximation therefore becomes
%%
\begin{align}
  \frac{dF_\lambda^\uparrow}{d\tau_{\lambda,z}} = D^\uparrow_\lambda \left( F_\lambda^\uparrow - \pi B_\lambda \right), \\
  \frac{dF_\lambda^\downarrow}{d\tau_{\lambda,z}} = -D^\downarrow_\lambda \left( F_\lambda^\downarrow - \pi B_\lambda \right).
\label{eq:rte_general_form_plane_parallel_flux_no_scattering}
\end{align}


The following sections describe the radiation schemes currently available in SCALE-DG.
Although the parameterizations of optical properties differ among the schemes, 
all are based on Eq.\,\eqref{eq:rte_general_form_plane_parallel_flux_no_scattering}.


\subsection{Gray radiation scheme}\label{sec:gray_radiation}
%%
As the simplest radiation model incorporating moisture feedback, a gray radiation model is provided. 
This scheme is gray in infrared so that a single optical thickness is defined for the entire longwave spectrum, 
which includes a parameterization of long-wave absorption by CO$_2$.


Following Eq.\,(5) in \cite{Vallis2018}, 
the vertical gradient of longwave optical thickness is parameterized as
%%
\begin{equation}
\frac{\partial \tau_{\mathrm{LW}}}{\partial \sigma}
=
a_{\mathrm{LW}}\mu + b_{\mathrm{LW}} q_v + c_{\mathrm{LW}} \ln\!\left(\frac{C_{\mathrm{CO_2}}}{360}\right),
\label{eq:gray_radiation_tau}
\end{equation}
%%
where $\sigma=p/P_0$, $q_v$ is the specific humidity,
$C_{\mathrm{CO_2}}$ is the atmospheric CO$_2$ concentration [ppm],
$\mu$ is a prescribed scaling factor,
and $a_{\mathrm{LW}}$, $b_{\mathrm{LW}}$, and $c_{\mathrm{LW}}$ are empirical coefficients.

\cite{Byrne2013} used $(a_{\mathrm{LW}}, b_{\mathrm{LW}}, c_{\mathrm{LW}})= (0.8678,\,1997.9,\,0)$
with the coefficients obtained by fitting this parameterization to the longwave optical depths employed by ~\cite{Frierson2006}.
Their idealized experiments used a planetary albedo of approximately $0.38$.
For experiments employing an albedo closer to that of Earth, approximately $0.30$, 
\cite{Vallis2018} suggested $(a_{\mathrm{LW}}, b_{\mathrm{LW}}, c_{\mathrm{LW}})= (0.1627,\,1997.9,\,0.17)$.
The latter parameter set also includes the explicit parameterization of longwave absorption by CO$_2$,
which was derived using output from the Santa Barbara DISORT Atmospheric Radiative Transfer model (SBDART)~\citep{Ricchiazzi1998}.


SCALE-DG uses the coefficient set of \cite{Vallis2018} by default, 
while the \cite{Byrne2013} set is also available as an alternative.


\subsection{Two-band infrared and one-band solar radiation scheme}\label{sec:two_band_radiation}
%%
This scheme has two infrared bands and one shortwave band, as described in \cite{Geen2016}, 
which provides an intermediate complexity between gray radiation and more sophisticated radiation schemes.
All bands were originally parameterized by fitting to data from SBDART for a range of atmospheric profiles. 
 \cite{Vallis2018} includes the CO$_2$ absorption in each band and changes the functional form of the non-window optical depth.
 The upward flux of shortwave radiation is assumed to be transparent.

Following Eq.\,(6) in \cite{Vallis2018},
the optical thickness in the shortwave band ($< 4~\mu\mathrm{m}$) is parameterized as
%%
\begin{equation}
\frac{\partial \tau_{\mathrm{SW}}}{\partial \sigma}
= a_{\mathrm{SW}} + b_{\mathrm{SW}}(\tau_{\mathrm{SW}}) q_v
+ c_{\mathrm{SW}} \ln\!\left(\frac{C_{\mathrm{CO_2}}}{360}\right),
\label{eq:2bandLW_1bandSW_rad_tau_sw}
\end{equation}
%%
where the humidity coefficient depends on the accumulated optical thickness as,
%%
\begin{equation}
\ln b_{\mathrm{SW}} (\tau_{\mathrm{SW}}) = \frac{0.01887}{\tau_{\mathrm{SW}}+0.009522}
+ \frac{1.603}{(\tau_{\mathrm{SW}}+0.5194)^2}.
\label{eq:twoband_bsw}
\end{equation}
%%
This dependence approximately represents the saturation of water-vapor absorption in the shortwave spectrum.

The longwave optical thickness is divided into a non-window band in 8-14$~\mu\mathrm{m}$ and an atmospheric window band in other long-wave wavelengths ($> 4~\mu\mathrm{m}$).
Following Eq.\,(7) in \cite{Vallis2018}, 
the non-window component is parameterized as
%%
\begin{equation}
\frac{\partial \tau_{\mathrm{LW}}^{\mathrm{nw}}}{\partial \sigma}
=
a_{\mathrm{LW}} + b_{\mathrm{LW}} \ln\!\left(1+c_{\mathrm{LW}} q_v\right)
+ d_{\mathrm{LW}} \ln\!\left(\frac{C_{\mathrm{CO_2}}}{360~\mathrm{ppm}}\right),
\label{eq:twoband_tau_lw}
\end{equation}
%%
while the atmospheric window band is expressed as
%%
\begin{equation}
\frac{\partial \tau_{\mathrm{LW}}^{\mathrm{win}}}{\partial \sigma}
= a_{\mathrm{win}} + b_{\mathrm{win}} q_v + c_{\mathrm{win}} q_v^2
+ d_{\mathrm{win}} \ln\!\left(\frac{C_{\mathrm{CO_2}}}{360~\mathrm{ppm}}\right).
\label{eq:twoband_tau_window}
\end{equation}
%%
The total longwave flux is obtained as the weighted sum of the two bands. 
Fixed weighting factors $R_{\mathrm{nw}}$ and $R_{\mathrm{win}}$ are prescribed for the non-window and window components, respectively.

The default parameter values adopted in the Isca implementation of Vallis et al.~\cite{Vallis2018} are 
%%
\begin{align}
(a_{\mathrm{SW}}, c_{\mathrm{SW}}) &= (0.0596,\,0.0029), \\
(a_{\mathrm{LW}}, b_{\mathrm{LW}}, c_{\mathrm{LW}}, d_{\mathrm{LW}}) &= (0.1000,\,23.8,\,254.0,\,0.2023), \\
(a_{\mathrm{win}}, b_{\mathrm{win}}, c_{\mathrm{win}}, d_{\mathrm{win}}) &= (0.2150,\,147.11,\,10814.0,\,0.0954).
\end{align}
%%
The partitioning factors for the longwave radiation are
$R_{\mathrm{nw}}=0.6268$ and $R_{\mathrm{win}}=0.3732$.


\subsection{Common formulation}
%%
The gray, one-band, and two-band radiation schemes share the same numerical procedure for evaluating radiative fluxes.
The parameterizations described in the preceding sections determine the optical-thickness gradient $\partial\tau/\partial\sigma$,
whereas the radiative transfer equations are solved using a common discretization and integration procedure.

Radiative transfer is solved independently for each vertical column.
The vertical grid used for the radiation calculation coincides with the DG nodal grid.
For each interval between adjacent DG nodes, the optical thickness for each band is
approximated by
%%
\begin{equation}
\Delta\tau = \left(\frac{\partial\tau}{\partial\sigma}\right) \Delta\sigma,
\label{eq:dtau}
\end{equation}
%%
where
$\Delta\sigma=\Delta p/P_0$,
and the optical-thickness gradient is evaluated using the arithmetic mean of the thermodynamic variables at the two adjacent nodes.
For the one-band solar radiation scheme,
the coefficient $b_{\mathrm{SW}}$ depends on the accumulated shortwave optical thickness and is therefore evaluated sequentially from the model top downward.

Radiative fluxes are calculated using the two-stream approximation without scattering.
Assuming that the optical thickness and the Planck source function are constant within each vertical interval,
the analytical solution of the radiative-transfer equation gives the upward and downward longwave fluxes as
%%
\begin{align}
F^{\uparrow}_{k+1} &= F^{\uparrow}_{k} e^{-D\Delta\tau}
+ \mathcal{B}\left(1-e^{-D\Delta\tau}\right), \\
F^{\downarrow}_{k} &= F^{\downarrow}_{k+1} e^{-D\Delta\tau} +\mathcal{B} \left(1-e^{-D\Delta\tau}\right),
\end{align}
%%
where $k$ is the vertical index, 
$D$ is the diffusivity factor, 
and $\mathcal{B}=\sigma_{\mathrm{SB}}T^4$ is the blackbody emission.
When the radiation schemes described in Secs.\,\ref{sec:gray_radiation} and \ref{sec:two_band_radiation} are used, we set $D=1$.
For the two-band scheme, the above calculation is applied independently to the non-window and window bands, 
and the total longwave flux is obtained as their weighted sum.

The downward shortwave flux is evaluated only by atmospheric attenuation,
%%
\begin{equation}
F^{\downarrow}_{k} = F^{\downarrow}_{k+1} e^{-\Delta\tau},
\end{equation}
%%
whereas the upward shortwave flux at the surface is determined by the prescribed surface albedo.
In the present implementation, the atmosphere is assumed to be
transparent to upward shortwave radiation.


Finally, the net radiative flux,
%%
\begin{equation}
F_{\mathrm{net}} = F^{\uparrow} - F^{\downarrow},
\end{equation}
%%
is differentiated in the vertical direction using the DG derivative operator to obtain the radiative heating rate.